\documentclass[../../main]{subfiles}

\begin{document}

\section{Conjunto das Partes}
\label{sec:matematica:teoria-conjuntos:conjunto-das-partes}

\begin{defn}[Conjunto das Partes]
	\label{def:matematica:teoria-conjuntos:conjunto-das-partes:conjunto-das-partes}

	Seja (A) um conjunto.

	O conjunto das partes de (A), denotado por

		[
			\mathcal{P}(A),
		]

	é o conjunto formado por todos os subconjuntos de (A).

	Formalmente,

	[
			\mathcal{P}(A)
				=
				{
					X
					\mid
					X\subseteq A
				}.
		]

\end{defn}

\begin{prop}[Caracterização do Conjunto das Partes]
	\label{prop:matematica:teoria-conjuntos:conjunto-das-partes:caracterizacao}

	Para quaisquer conjuntos (A) e (X),

	[
			X\in\mathcal{P}(A)
			\Leftrightarrow
			X\subseteq A.
		]

\end{prop}

\begin{proof}

	A equivalência segue diretamente da definição.

\end{proof}

\begin{ex}

	Seja

	[
	A={1,2}.
	]

	Então

	[
	\mathcal{P}(A)
		=
		{
			\varnothing,
			{1},
			{2},
			{1,2}
		}.
	]

\end{ex}

\begin{prop}
	\label{prop:matematica:teoria-conjuntos:conjunto-das-partes:vazio}

	Para todo conjunto (A),

	[
			\varnothing
			\in
			\mathcal{P}(A).
		]

\end{prop}

\begin{proof}

	Como

	[
	\varnothing\subseteq A,
	]

	segue pela definição que

		[
			\varnothing\in\mathcal{P}(A).
		]

\end{proof}

\begin{prop}
	\label{prop:matematica:teoria-conjuntos:conjunto-das-partes:proprio}

	Para todo conjunto (A),

	[
			A
			\in
			\mathcal{P}(A).
		]

\end{prop}

\begin{proof}

	Todo conjunto é subconjunto de si próprio.

	Portanto

		[
			A\subseteq A.
		]

	Logo

		[
			A\in\mathcal{P}(A).
		]

\end{proof}

\begin{prop}[Monotonicidade]
	\label{prop:matematica:teoria-conjuntos:conjunto-das-partes:monotonicidade}

	Se

		[
			A\subseteq B,
		]

	então

	[
	\mathcal{P}(A)
	\subseteq
	\mathcal{P}(B).
	]

\end{prop}

\begin{proof}

	Seja

	[
	X\in\mathcal{P}(A).
	]

	Então

	[
	X\subseteq A.
	]

	Como

		[
			A\subseteq B,
		]

	segue que

		[
			X\subseteq B.
		]

	Portanto

		[
			X\in\mathcal{P}(B).
		]

\end{proof}

\begin{prop}[Caracterização da Inclusão]
	\label{prop:matematica:teoria-conjuntos:conjunto-das-partes:inclusao}

	Para quaisquer conjuntos (A) e (B),

	[
			A\subseteq B
			\Leftrightarrow
			\mathcal{P}(A)
			\subseteq
			\mathcal{P}(B).
		]

\end{prop}

\begin{proof}

	Já foi demonstrada a implicação direta.

	Para a recíproca, suponha

		[
			\mathcal{P}(A)
			\subseteq
			\mathcal{P}(B).
		]

	Como

		[
			A\in\mathcal{P}(A),
		]

	segue que

		[
			A\in\mathcal{P}(B).
		]

	Portanto

		[
			A\subseteq B.
		]

\end{proof}

\end{document}
